\documentclass[11pt]{letter}

\usepackage{geometry}
\geometry{margin=1in}
\usepackage{setspace}
\usepackage{hyperref}

\signature{Decory Edwards}
\address{Decory Edwards \\ Johns Hopkins University \\ Department of Economics}

\begin{document}

\begin{letter}{Hiring Committee \\ Private Equity Analyst Program \\ New York, NY}

\opening{Dear Hiring Committee,}

I am writing to apply for the 2026 Private Equity Analyst position in New York. I am currently a PhD candidate in Economics at Johns Hopkins University. My training combines rigorous quantitative analysis, macroeconomic reasoning, and disciplined evaluation of risk and uncertainty. I am motivated by work that requires careful diligence, clear judgment, and the integration of qualitative and quantitative information to support investment decisions. The breadth of exposure across manager diligence, co-investments, and secondaries offered by this role aligns closely with my interests and analytical approach.

My research agenda is at the intersection of wealth inequality and macroeconomics. In my job market paper, I build and estimate a structural heterogeneous agent model with returns that differ across households. The core contribution is to show how heterogeneity in returns and income risk jointly shape the wealth distribution and consumption behavior. Relative to closely related work, my framework performs better at matching the lower and middle portions of the wealth distribution. This difference matters quantitatively. Models that focus primarily on returns heterogeneity can fit the top of the wealth distribution closely but tend to generate too much wealth among the bottom half of households. In my framework, income uncertainty generates a precautionary saving motive that produces genuinely low wealth households. This leads to a realistic distribution of marginal propensities to consume and aggregate consumption responses that align with empirical estimates.

A key feature of my research is the disciplined assessment of economic environments and risk–return tradeoffs. I work with large datasets, analyze macroeconomic and sector-level trends, and translate findings into structured written analysis. I am comfortable evaluating assumptions, stress testing outcomes, and synthesizing information from multiple sources into concise recommendations. This work maps directly to manager and asset-level diligence, market research, and the preparation of investment committee materials.

I am particularly interested in this role because of its emphasis on fundamental analysis and its exposure across private equity strategies and geographies. The opportunity to engage directly with fund managers and portfolio companies, conduct reference calls, and contribute to underwriting and monitoring decisions is especially appealing. I value environments that emphasize teamwork, intellectual rigor, and sound investment judgment.

My economics training has provided a strong foundation in quantitative reasoning, financial analysis, and clear communication. I am highly proficient in Excel, PowerPoint, and Word and take pride in producing well-structured analyses with careful attention to detail. I work effectively in team settings, manage multiple workstreams, and communicate ideas clearly in both written and verbal form.

I am enthusiastic about the opportunity to begin my career in private equity through this analyst program in New York. I see this role as a strong platform to apply my analytical training to long-term investing while developing a broad understanding of private markets.

I would welcome the opportunity to contribute my skills and perspective to your investment team. Thank you for your time and consideration.

\closing{Sincerely,}

\end{letter}
\end{document}
